\Needspace{130bp}
\section{交通标志图像分类与参数影响}

\subsection{分类结果与特征维度}
将数据集更换为10类交通标志，使用6,024张训练图像和324张测试图像。
图像保持 $224\times224$ 原尺寸，灰度化后提取HOG特征，使用线性SVM分类，固定 $C=1$。
已尝试配置中，方向数9、Cell大小 $4\times4$ 像素、Block大小 $4\times4$ 个Cell的结果最好，
测试准确率为95.06\%（308/324），比原始参数提高1.54个百分点，测试正确数增加5张。
此时图像包含 $56\times56$ 个Cell，可放置 $53\times53$ 个Block，特征维度为
$D=53\times53\times4\times4\times9=404\,496$。

\ReportTableCaption{交通标志分类结果比较\label{tab:traffic-classification}}
\input{traffic-classification-table.tex}
HOG参数顺序与表\ref{tab:cifar-classification}相同；各组训练准确率均为100\%。
测试集参与了HOG参数比较，结果限于本次候选配置及该数据划分。

\subsection{HOG参数影响分析}
\textbf{Cell与Block大小。}固定方向数9、Block为 $2\times2$ 时，
Cell边长4像素的测试准确率为94.14\%，高于8和16像素的93.52\%，相差2张图像。
在Cell为 $4\times4$ 时，将Block由 $2\times2$ 增大至 $4\times4$，准确率进一步提高到95.06\%。
较小Cell与较大Block的配合在本次比较中较好，但特征维度也明显增加。

\textbf{方向数。}固定Cell为 $4\times4$、Block为 $4\times4$ 时，方向数从6增加到9，
准确率由93.52\%提高到95.06\%；继续增加到12反而降至94.14\%。
这说明方向划分更细并不一定带来更好的分类效果，本次三组中9个方向较合适。

\subsection{分类样本对比}
\begin{minipage}{\linewidth}
  \begin{minipage}[c]{0.49\linewidth}
    \begin{minipage}[c]{72bp}\includegraphics[width=\linewidth]{traffic-correct.png}\end{minipage}\hfill
    \begin{minipage}[c]{0.67\linewidth}
      \textbf{分类正确}\\真实类别：警告标志\\模型预测：警告标志
    \end{minipage}
  \end{minipage}\hfill
  \begin{minipage}[c]{0.49\linewidth}
    \begin{minipage}[c]{72bp}\includegraphics[width=\linewidth]{traffic-incorrect.png}\end{minipage}\hfill
    \begin{minipage}[c]{0.67\linewidth}
      \textbf{分类错误}\\真实类别：警告标志\\模型预测：人行横道标志
    \end{minipage}
  \end{minipage}
  \captionof{figure}{交通标志测试样本及中文预测类别，使用HOG 9/4/4、$C=1$ 的模型。}
  \label{fig:traffic-examples}
\end{minipage}
